\section{数学归纳法为何成立}
\label{sec:01-Mathematical-Language/04-Induction-and-Recursion/01}

手算几个和：

\[
\begin{aligned}
1 &= 1 = 1^2,\\
1+3 &= 4 = 2^2,\\
1+3+5 &= 9 = 3^2,\\
1+3+5+7 &= 16 = 4^2.
\end{aligned}
\]

前 \(n\) 个正奇数相加，看上去总是 \(n^2\)。再往下试，\(n=10\) 给出 100；\(n=100\) 用等差数列公式一验，也对。信心在涨。

可是把 \(n\) 推到 \(10^{100}\) 呢？你手上这张验证表，无论加长到多少行，后面永远还剩无穷多个没检查的 \(n\)。这不是耐心问题，也不是算力问题：有限长的验证表和``对所有自然数成立''之间，隔着一条靠列举跨不过去的沟。在\hyperref[ch:006]{``证明：让结论从已知条件中不得不成立''}一章里已经说过，例子再多也只是证据。

\subsection{自然数是造出来的，不是摆在那里的}

跨过这条沟靠的不是更快的验证，而是换一个角度看自然数本身。

\(7\) 不是凭空躺在数轴上的一个符号，它是 \(0\) 连续执行七次``加一''的产物。\(10^{100}\) 也一样，只是步数多些。全体自然数不是一堆散装对象，而是一个起点加一条生成规则反复作用的结果——每个自然数都有一条从 \(0\) 通到它的、长度有限的生成路径。

这个观察改变了任务的形状。既然对象是被规则造出来的，就不必逐个检查对象，只要盯住规则。

\begin{quote}
用有限的论证覆盖无限的对象，靠的是对象自身的生成结构：论证跟着生成规则走，被生成的一切就都被覆盖了。
\end{quote}

自然数只有一条生成规则，所以论证也只需要两件东西：起点上命题为真，以及那条规则不破坏真。

\begin{center}
\begin{tikzpicture}[x=1cm,y=1cm,>=stealth,font=\small]
  \foreach \x/\lab in {0/0, 1.2/1, 2.4/2, 3.6/3, 5.4/k, 6.6/{k+1}}{
    \draw[fill=black!8] (\x,0) rectangle (\x+0.3,1.0);
    \node[below] at (\x+0.15,-0.05) {\(\lab\)};
  }
  \foreach \x in {0, 1.2, 2.4}{
    \draw[->,black!60] (\x+0.42,0.5) -- (\x+1.08,0.5);
  }
  \draw[->,black!60] (5.82,0.5) -- (6.48,0.5);
  \node at (4.6,0.5) {\(\cdots\)};
  \draw[->,thick] (-1.15,0.5) -- (-0.15,0.5);
  \node[above] at (-0.65,0.55) {基例};
  \node at (2.9,1.75) {归纳步：每一处都是同一段论证};
\end{tikzpicture}

\emph{基例把第一张牌推倒，归纳步保证任意相邻两张之间的传递；两者合起来，链条上没有一张牌能幸免。}
\end{center}

\begin{theorem}
设 \(P(n)\) 是关于整数 \(n\ge n_0\) 的命题。若 \(P(n_0)\) 成立，且对每个 \(k\ge n_0\) 都能由 \(P(k)\) 推出 \(P(k+1)\)，则 \(P(n)\) 对所有 \(n\ge n_0\) 成立。
\end{theorem}

归纳假设的措辞里藏着一处精细的区分：你假设的是``某个任意但固定的 \(k\) 满足 \(P(k)\)''，而不是``所有 \(n\) 都满足 \(P(n)\)''。前者是局部的、临时的，只为证明从 \(k\) 到 \(k+1\) 这一步的传递；后者是最终目标。把两者混为一谈，才会变成循环论证。

\subsection{把猜想钉成定理}

回到开头的和式。命题 \(P(n)\)：\(1+3+\cdots+(2n-1)=n^2\)，\(n\ge 1\)。

基例 \(P(1)\)：左边只有 \(1\)，右边是 \(1^2\)。

归纳步：设 \(P(k)\) 成立。第 \(k+1\) 个奇数是 \(2k+1\)，于是

\[
\underbrace{1+3+\cdots+(2k-1)}_{\text{由 }P(k)\text{ 等于 }k^2}+(2k+1)
= k^2+2k+1=(k+1)^2 ,
\]

这正是 \(P(k+1)\) 的结论。

代数只有一行。真正吃力的是它前面那一步：把 \(P(k+1)\) 的左边拆成``\(P(k)\) 的左边''加``新增的一项''。归纳步能不能走通，几乎总是取决于能不能在规模 \(k+1\) 的对象里认出规模 \(k\) 的对象。认不出这个拆解，代数配得再平也没用。

\subsection{两个条件各自防住什么}

基例和归纳步不是仪式感的两段式，它们各自封住一种失效。

先看没有基例会怎样。考虑 \(Q(n)\)：\(1+2+\cdots+n=\frac{n(n+1)}{2}+7\)。归纳步无懈可击：设 \(Q(k)\) 为真，两边同加 \(k+1\)，得到

\[
\frac{k(k+1)}{2}+7+(k+1)=\frac{(k+1)(k+2)}{2}+7 ,
\]

正是 \(Q(k+1)\)。而 \(Q(1)\) 说的是 \(1=8\)，假；不仅如此，\(Q(n)\) 对每一个 \(n\) 都假。传递机制装配完好，链条一节不缺，只是整条链没有一端接在真命题上。基例的职责就是这个接头。

反过来，基例成立而归纳步在某个 \(k\) 上悄悄失效，同样出事。经典的例子是``任意 \(n\) 匹马颜色相同''：\(n=1\) 显然成立；归纳步里，把 \(k+1\) 匹马编号，去掉最后一匹得到前 \(k\) 匹同色，去掉第一匹得到后 \(k\) 匹同色，两组有重叠的马，于是全体同色。

论证读着顺，结论却荒谬。断裂点在 \(k=1\)：此时两组各只剩一匹马，中间没有共同的马来传递颜色。归纳步用到的那个结构——``两组有重叠''——在最小的 \(k\) 上根本不存在。归纳步必须对每一个 \(k\) 成立，写完之后回头验一遍最小规模，是很便宜的保险。

同类错误还有一种更隐蔽的形态：归纳步里不知不觉调用了 \(P(k-1)\)。手里只有 \(P(k)\)，却用了两个历史情形，链条同样断了。下一节专门处理这种依赖。

\subsection{起点和步长划定覆盖范围}

归纳法覆盖到哪里，由基例的位置和归纳步的跨度共同决定。命题若从 \(n=5\) 才开始成立，基例就得放在 \(P(5)\)；只验证 \(P(1)\) 跨不过中间那段断层。

跨度也一样。假如归纳步只能从 \(P(k)\) 推到 \(P(k+2)\)，一个基例就只能覆盖一条同奇偶的链。

\begin{center}
\begin{tikzpicture}[x=1cm,y=1cm,>=stealth,font=\small]
  \foreach \i in {0,...,6}{
    \node at (\i*1.2,0) {\(\i\)};
  }
  \foreach \i in {0,2,4}{
    \draw[->,thick] (\i*1.2+0.22,0.28) .. controls (\i*1.2+0.5,1.0) and (\i*1.2+1.9,1.0) .. (\i*1.2+2.18,0.28);
  }
  \foreach \i in {1,3,5}{
    \node[red!65!black] at (\i*1.2,-0.7) {\(\times\)};
  }
  \draw[->,thick] (-1.1,0) -- (-0.4,0);
  \node[above] at (-0.75,0.05) {基例};
  \node[right] at (7.6,0) {奇数一个也没被覆盖};
\end{tikzpicture}

\emph{步长为 2 时，一个基例只能把传递力送到同奇偶的那条链上；要覆盖全体，还得补一个基例 \(P(1)\)。}
\end{center}

所以下笔之前先问一句：我要让这股传递力抵达哪些对象？再据此挑起点、定步长，让传递路径真的经过每一个目标。

\subsection{它负责验证，不负责发现}

归纳法把一条已经写清楚的猜想推到全体自然数，但它不会告诉你猜想从哪来。前 \(n\) 个奇数之和为什么猜 \(n^2\)？可能来自用石子摆正方形的几何直觉，可能来自算了十几项后看出的数值规律，也可能来自解一个递推式。这些都属于发现阶段，归纳法一概不管。

把发现和验证分开，两件工具才各就各位：数值实验负责提出候选，几何或代数的洞察负责解释为什么值得相信，归纳法负责最后一跃——把有限的验证拉过全称量词那条沟。

下一节要拿走一个限制。这里的归纳步只准使用 \(P(k)\)，而 Fibonacci 数列、分治算法、素数分解这类对象的下一步要回望的不止一个历史情形，归纳假设得开得更大。
